\begin{lemma}[Sequential completion of a colour class]
Let $B$ be a blow-up instance with finite source and host vertex sets $V_H,V_G$,
candidacy system $A$, and reduced graph $R$ on $[r]$.  Fix a colouring
$c:[r]\to[q]$, sequences $\varepsilon,p:\mathbb N\to[0,\infty]$, a degree
bound $D\in\mathbb N$, $C\in[0,\infty]$, and $n,t\in\mathbb N$.
Write $\mathcal E(Z)$ for the compatible partial embeddings with domain $Z$,
$B_t=\{j:c(j)=t\}$, and
\[
 X(t,T)=\{x\in V_H:c(B.\mathrm{sourceCluster}(x))<t
                     \text{ or }B.\mathrm{sourceCluster}(x)\in T\}.
\]
Use the invariant $\mathcal I(t,T;Z,\nu)$ of
\noderef{ColourRoundInvariant}, with parameters $(\varepsilon,p,D,C,n)$.
Assume the following one-cluster transition for every finite $T\subseteq[r]$,
finite $Z\subseteq V_H$, and specified law $\nu$ on $\mathcal E(Z)$:
if $T\subseteq B_t$ and $\mathcal I(t,T;Z,\nu)$ holds, then for every
$i\in B_t\setminus T$, setting $Z'=X(t,T\cup\{i\})$ and
$Y_i=\{x:B.\mathrm{sourceCluster}(x)=i\}$, there is a probability kernel
$\kappa:\mathcal E(Z)\to\operatorname{PMF}(\mathcal E(Z'))$ such that
\begin{enumerate}
\item for every $f\in\operatorname{supp}\nu$, every
$h\in\operatorname{supp}\kappa(f)$, and every $x\in Z$, $h(x)=f(x)$;
\item for each $f\in\operatorname{supp}\nu$, the restriction to $Y_i$ of the
law $\kappa(f)$ is $(C/n)$-vertex-spread on the entire host vertex set;
\item $\mathcal I(t,T\cup\{i\};Z',\nu\mathbin{\mathrm{bind}}\kappa)$ holds
for the actual bound law.
\end{enumerate}
Then, for every $S\subseteq B_t$, finite $X\subseteq V_H$, and specified
law $\mu$ on $\mathcal E(X)$ satisfying $\mathcal I(t,S;X,\mu)$,
there is a kernel
$K:\mathcal E(X)\to\operatorname{PMF}(\mathcal E(X(t,B_t)))$ such that
every supported output of $K(f)$ agrees with $f$ on $X$ whenever
$f\in\operatorname{supp}\mu$, and
\[
 \mathcal I(t,B_t;X(t,B_t),\mu\mathbin{\mathrm{bind}}K).
\]
In particular this completed law is $(C/n)$-vertex-spread.
\end{lemma}
\begin{proof}
This is the finite sequential control argument in
\texttt{paper/r2-colour-round-addendum.tex}, lines 241--271 and 333--341.
All invariants in this proof have the same round $t$; no new round is created.
By \noderef{ColourRoundInvariant}, the starting domain is exactly $X(t,S)$,
and the law's spread clause is the restriction-law bound of
\noderef{VertexSpreadOn}.  Maps and their restrictions are interpreted as in
\noderef{PartialCompatibleEmbedding}.

Induct on the finite remaining set $F=B_t\setminus S$, equivalently on
$|F|$, with the induction assertion quantified over all $S,X,\mu$ satisfying
the stated assumptions.  If $F$ is empty, $S=B_t$ since $S\subseteq B_t$.
The domain clause gives $X=X(t,B_t)$.  Identify the two partial-embedding
types using this equality, without changing any map.  Let $K(f)$ be the
point mass at $f$ under this identification.  Binding $\mu$ to these point
masses is precisely $\mu$ under the same domain identification, so its
invariant is the given one and map preservation is immediate.  This also
covers an empty colour class.  No transition is invoked in this case.

If $F$ is nonempty, choose $i\in F$.  Then $c(i)=t$ and $i\notin S$,
so apply the stated transition to the actual current law $\mu$.  Retain
its kernel $\kappa$, put $S_1=S\cup\{i\}$, $X_1=X(t,S_1)$, and put
$\mu_1=\mu\mathbin{\mathrm{bind}}\kappa$.
The transition supplies $\mathcal I(t,S_1;X_1,\mu_1)$, as well as preservation
of $X$ and the conditional spread of the new cluster for every supported
past outcome.  Since $i\in B_t\setminus S$, $S_1\subseteq B_t$ and
$B_t\setminus S_1=(B_t\setminus S)\setminus\{i\}$ has cardinality
$|F|-1$.  The induction hypothesis applied to this specified $\mu_1$
therefore gives a kernel $L$ from $\mathcal E(X_1)$ to
$\mathcal E(X(t,B_t))$, preserving every supported input on $X_1$, and
\[
 \mathcal I(t,B_t;X(t,B_t),\mu_1\mathbin{\mathrm{bind}}L).
\]
Define $K(f)=\kappa(f)\mathbin{\mathrm{bind}}L$ for every input $f$, including
inputs outside $\operatorname{supp}\mu$; no properties are required there.
The probability law of this composite kernel satisfies
\[
 \mu\mathbin{\mathrm{bind}}K
   =(\mu\mathbin{\mathrm{bind}}\kappa)\mathbin{\mathrm{bind}}L
   =\mu_1\mathbin{\mathrm{bind}}L.
\]
Indeed, the mass of an output $h$ on either side is
$\sum_f\sum_g\mu(f)\kappa(f)(g)L(g)(h)$; reassociation of the nonnegative
probability sums proves the equality.  Thus the final invariant is for
the recursively bound law, not for a newly chosen unrelated law.

To check preservation, fix $f\in\operatorname{supp}\mu$ and
$h\in\operatorname{supp}K(f)$.  The bind mass formula gives some
$g\in\operatorname{supp}\kappa(f)$ with $h\in\operatorname{supp}L(g)$:
a positive nonnegative sum has a positive summand, and a positive product
has positive factors.  The same formula, using the positive masses
$\mu(f)$ and $\kappa(f)(g)$, shows $g\in\operatorname{supp}\mu_1$.
Also $X\subseteq X_1$, because $X=X(t,S)$ and $S\subseteq S_1$.
For $x\in X$, preservation for $L$ and then for $\kappa$ gives
$h(x)=g(x)=f(x)$.  This proves the asserted preservation for $K$ and completes
the finite induction.  At each stage the spread clause is supplied for that
stage's bound law by the transition and retained by the induction hypothesis;
unfolding \noderef{ColourRoundInvariant} gives the final spread conclusion.
The conditional new-cluster bounds are those of the successive kernels
conditional on their actual past outcomes.  No conditioning on a global
success event, and no multiplication by a class-size-dependent loss, occurs.
\end{proof}
